\subsection{单环的理想}

设 D 是一个体，V 是 D 上一个维数有限 \(n \geqslant 1\) 的右向量空间。考虑单环 \(A = \text{End}_{\text{D}}(\text{V})\)。对 V 的每个线性子空间 W，把 A 中满足 \(aW = 0\)（相应地，\(aV \subset W\)）的元素 a 之集记作 \(\mathfrak{a}(\text{W})\)（相应地，\(\mathfrak{b}(\text{W})\)）。

命题 7. —— a) 映射 \(W \mapsto \mathfrak{a}(W)\) 是从 V 的线性子空间之集到 A 的左理想之集的双射。

\begin{enumerate}
\def\labelenumi{\alph{enumi})}
\setcounter{enumi}{1}
\item
  映射 \(W \mapsto \mathfrak{b}(W)\) 是从 V 的线性子空间之集到 A 的右理想之集的双射。
\item
  设 \(\mathrm{W}_1\) 与 \(\mathrm{W}_2\) 是 \(\mathrm{V}\) 的线性子空间。关系 \(\mathrm{W}_1 \subset \mathrm{W}_2\)、\(\mathfrak{a}(\mathrm{W}_1) \supset \mathfrak{a}(\mathrm{W}_2)\) 与 \(\mathfrak{b}(\mathrm{W}_1) \subset \mathfrak{b}(\mathrm{W}_2)\) 等价。
\end{enumerate}

断言 b) 由把 VIII, 第 104 页的例 1, b) 应用于可逆 \((\mathrm{D}^{\circ},\mathrm{A}^{\circ})\)-双模 V 得出，关系 \(\mathrm{W}_1\subset \mathrm{W}_2\) 与 \(\mathfrak{b}(\mathrm{W}_1)\subset \mathfrak{b}(\mathrm{W}_2)\) 的等价性亦然。

设 \(V^{*}\) 是 V 的对偶，视为 D 的反体 \(D^{\circ}\) 上的右向量空间。对 V 的每个子空间 W，把 W 在 \(V^{*}\) 中的正交子空间记作 \(W'\)。映射 \(W \mapsto W'\) 是从 V 的子空间之集到 \(V^{*}\) 的子空间之集的双射。若 \(W_{1}\) 与 \(W_{2}\) 是 V 的两个子空间，则关系 \(W_{1} \subset W_{2}\) 与 \(W_{1}' \supset W_{2}'\) 等价。而映射 \(u \mapsto {}^{t}u\) 是从 A 到 \(\operatorname{End}_{D^{\circ}}(V^{*})\) 的反环的同构；它把 A 的左理想变为 \(\operatorname{End}_{D^{\circ}}(V^{*})\) 的右理想，并把 \(\mathfrak{a}(W)\) 变为 \(\mathfrak{b}(W')\)，即 \(V^{*}\) 的满足 \(h(V^{*}) \subset W'\) 的自同态 h 之集。于是断言 a) 以及关系 \(W_{1} \subset W_{2}\) 与 \(\mathfrak{a}(W_{1}) \supset \mathfrak{a}(W_{2})\) 的等价性，由关于 V 的对偶 \(V^{*}\) 的与 b) 类似的断言得出。

推论. —— a) A 的极小左理想是诸理想 \(\mathfrak{a}(\mathrm{H})\)，其中 H 是 V 的一个超平面。A 的极大左理想是诸理想 \(\mathfrak{a}(\mathrm{L})\)，其中 L 是 V 的一条直线。

\begin{enumerate}
\def\labelenumi{\alph{enumi})}
\setcounter{enumi}{1}
\tightlist
\item
  A 的极小右理想是诸理想 \(\mathfrak{b}(\mathrm{L})\)，其中 L 是 V 的一条直线。A 的极大右理想是诸理想 \(\mathfrak{b}(\mathrm{H})\)，其中 H 是 V 的一个超平面。
\end{enumerate}

设 \((\mathrm{L}_{i})_{i\in\mathrm{I}}\) 是一族直和为 V 的直线。设 \((\varepsilon_{i})_{i\in\mathrm{I}}\) 是与分解 \(V=\oplus_{i\in\mathrm{I}}L_{i}\) 相伴的投影算子族。诸 \(\varepsilon_{i}\) 是 A 中的幂等元，并且我们有 \(\varepsilon_{i}\varepsilon_{j} = 0\)（对 \(i\neq j\)）且 \(\sum_{i\in \mathrm{I}}\varepsilon_i = 1\)。把超平面 \(\sum_{j\neq i}\mathrm{L}_j\) 记作 \(\mathrm{H}_i\)；它是 \(\varepsilon_{i}\) 的核。于是我们有

\[
\mathfrak {a} (\mathrm{H} _ {i}) = \mathrm{A}   \varepsilon_ {i}  , \quad \mathfrak {b} (\mathrm{L} _ {i}) = \varepsilon_ {i}   \mathrm{A}  .
\]

A-模 \(A_{s}\) 是极小左理想族 \((\mathfrak{a}(\mathrm{H}_{i}))_{i\in\mathrm{I}}\) 的直和，而 \(A_{d}\) 是极小右理想族 \((\mathfrak{b}(\mathrm{L}_{i}))_{i\in\mathrm{I}}\) 的直和。

考虑特殊情形 \(V = D_{d}^{n}\)，并把 A 等同于矩阵环 \(\mathbf{M}_{n}(\mathrm{D})\)。用 I 表示 N 中的区间 \([1, n]\)，用 \((v_{i})_{i \in \mathrm{I}}\) 表示 V 的典范基；令 \(L_{i} = v_{i}D\)，并用 \(E_{ij}\) 表示矩阵单位（II, §10, No.~3, 第 341 页）。于是我们有 \(\varepsilon_{i} = E_{ii}\)。左理想 \(AE_{ii}\) 等于 \(DE_{1i} + \cdots + DE_{ni}\)，由除第 i 列外所有列均为零的矩阵组成。右理想 \(E_{ii}A\) 等于 \(DE_{i1} + \cdots + DE_{in}\)，由除第 i 行外所有行均为零的矩阵组成。我们还有关系式

\[
\mathrm{E} _ {i i} \mathrm{AE} _ {j j} = \mathrm{E} _ {i i} \mathrm{A} \cap \mathrm{AE} _ {j j} = \mathrm{DE} _ {i j}
\]

其中 i 与 j 介于 1 与 n 之间。
% semantic-fragments: legacy-input-seam
\relax{}
